\section{Near-maximum ordinary small-tail layers}
\label{sec:near-maximum-small-tail}

\paragraph{Review status and scope.}
D-TD51, question Q2fg, is a new derivation by
\texttt{/root/imprimitive}.  Complete author, independent and root
whole-proof reviews pass; their exact scopes are recorded separately.
It uses the accepted fixed-tail local
estimates, the unshortened weighted-cycle inequality in HD-TW04,
and the induced-cover geometry isolated in D-TD49.  The complete
original coefficient source is retained throughout.  In particular
no rank-six budget or twelve projected marks are assigned to a
layer or to its external partner.

Put
\[
 A_*=240!\,2^{239},\quad M_*=1+241A_*,
 \quad F_*=240M_*+2,\quad D_*=100(1+M_*),
 \quad N_*=(D_*F_*^{240})^2.
 \tag{NMT.1}
\]
This is exactly the existing HD-ML01/HD-TW03 cutoff; it is not enlarged.

\begin{theorem}[D-TD51]
Let an actual connected finite $J$-Galois sphere cover have a
faithful primitive complex coefficient representation $V$ and
genuine generating product-one inertia $T_1,\ldots,T_r$ satisfying
\[
 e=\dim V,\qquad r\ge12,\qquad
             \sum_i\operatorname{rank}(1-T_i)=2e+6.
 \tag{NMT.2}
\]
Fix one complete D-TD12 decomposition, let $n\ge N_*$ be its
largest alternating natural degree, and suppose it contains an
actual transitive alternating extraction layer with
\[
 m\ge\left\lceil\frac{n+120}{108}\right\rceil,\qquad
 V=(W_1\otimes\cdots\otimes W_t)\otimes R,\quad
 \dim W_j=a,\quad t\ge1.
 \tag{NMT.3}
\]
Suppose its component module is the ordinary irreducible
$A_m$-restriction of $S^{(m-k,\mu)}$, $\mu\vdash k$, or the
transpose label, for some $2\le k\le240$.
Then this source is impossible.

The entire complement $R$ is retained.  Neither an independent
source for $R$ nor equality of coefficient fields is assumed.
Together with D-TD49, the conclusion says that a surviving
near-maximum layer in this range must be ordinary standard;
it does not yet say that the standard layer is a singleton.
\end{theorem}

\begin{proof}
The accepted all-layer support floor D-TD13/HD-FT21 gives, for
every nonscalar operator $B$ in the actual entire complement,
\[
 \frac{\nu_R(B)}{\dim R}\ge\frac1{n-1}\ge\frac1{108m},
 \qquad
 \nu(B)=\dim B-\max_z\dim\ker(B-zI).
 \tag{NMT.4}
\]
If the complement is one-dimensional this is vacuous, and genuine
total scalar entries have full actual moving rank.  The inequality
uses the same complete original recursion: all its remaining
alternating natural degrees are at most $n$.  The layer image
contains $A_m^t$ and is transitive on the components, so its
actions on $t$ components and $mt$ natural letters are actual
transitive quotient actions.  Compatible cycle products are
represented by the indicated ordinary $S_m$-module up to scalars.
Fixed component identifications and all scalar phases are handled
as in the accepted ordinary normalizer interface.

Since $N_*\ge10000\cdot241^{482}$, (NMT.3) implies
\[
                       m>241^{482}>10^{15}.
 \tag{NMT.5}
\]
This also ensures $m>2k+2$, so the displayed ordinary restrictions
are irreducible and not split.

\subsection*{1. Tightening the existing fixed-tail threshold}

For $3\le k\le240$, write the accepted stable character polynomial
as $q_\mu=CX_1^k+R_\mu$, with
$C=f^\mu/k!$ and
$M_\mu=1+C^{-1}\|R_\mu\|_1$.
The proof of (TW.23) already establishes
\[
 \|q_\mu\|_1\le f^\mu\sum_{j=0}^kp(j).
 \tag{NMT.6}
\]
Use the sum of the elementary bounds, rather than their maximum:
$p(0)=1$ and $p(j)\le2^{j-1}$ for $j\ge1$ give
$\sum_{j=0}^kp(j)\le2^k$.
Removing the leading monomial cannot increase coefficient norm.
Thus the following deliberately weaker inequality suffices:
\[
 M_\mu\le1+k!2^k\le1+2A_*<\frac{M_*}{100},
 \qquad
 F_\mu=\lceil kM_\mu\rceil+1\le\frac{F_*}{2}.
 \tag{NMT.7}
\]
The strict comparison follows from $41A_*>99$.
For the last one use $k\le240$, $M_\mu<M_*/100$, and
$2.4M_*+2<120M_*+1=F_*/2$.

Recall the full accepted threshold, not just its largest term:
\[
 N_\mu=\left\lceil\max\left\{
 2k+2,\,200,\,100/C,\,100kM_\mu,\,
 [100(1+M_\mu)]^5,\,
 [100(1+M_\mu)F_\mu^k]^2\right\}\right\rceil.
 \tag{NMT.8}
\]
All terms other than the last are at most $D_*^5$:
$C^{-1}\le240!<M_*$, $k\le240$, and $D_*>482$ suffice
also for the linear terms.  The last term is at most
$N_*/2^{480}$, by (NMT.7), $k\le240$, and $F_*/2>1$.
Moreover $F_*\ge D_*$, so
$N_*\ge D_*^{482}$ and $D_*^5\le N_*/108$.
As $2^{480}>108$,
\[
                    N_\mu\le\left\lceil N_*/108\right\rceil<m.
 \tag{NMT.9}
\]
This is an elementary refinement of the already proved coefficient
estimate.  No new stable-character formula or asymptotic constant
has been introduced.

The local all-target bounds (TW.26)--(TW.27) therefore hold.
For tails $k\ge3$ they give, with $S=Cm^{k-1}$,
\[
 \begin{gathered}
 .99mS\le a\le1.01mS,\qquad
 A=a/m-1.01S/m,\quad B_{\rm old}=.4S,\\
 A\ge0,\qquad
 \lambda=1-(mA+1)/a\ge0,\qquad
 1\le L=a-(m-1)A\le2.02S<5a/m,\\
 \nu(W(y))\ge A\operatorname{ind}(y)+B_{\rm old}
                       \quad(y\ne1).
 \end{gathered}
 \tag{NMT.10}
\]
For tail two use the accepted exact bounds instead:
$A=(m-5)/2$, $B_{\rm old}=m/4$,
$a=m(m-3)/2$ or $(m-1)(m-2)/2$.
Their two values of $L$ are $(3m-5)/2$ and $(3m-3)/2$;
again $\lambda\ge0$, $L\ge1$, and $L\le5a/m$.

In all these cases $a\ge m^2/4$.
For tail two this is immediate at $m\ge64$.
For $k\ge3$ it also follows directly, without an approximate
coefficient argument, from the accepted MSP Lemma~8(ii) tableau
bound:
\[
 a\ge\binom{m-k}{k}f^\mu
   \ge\left(\frac{m}{2k}\right)^k
   \ge(m/480)^k\ge m^2.
 \tag{NMT.11}
\]
Here $m\ge2k$ and $m\ge480^3$; the last inequality follows from
$3(k-2)\ge k$ for $k\ge3$.

\subsection*{2. At most 239 components and three component-moving marks}

Set
\[
                         B'=\frac{a}{108m}.
 \tag{NMT.12}
\]
The previous local $B_{\rm old}$ is at least $B'$ in every case.
For $k\ge3$, use $a\le1.01mS$; for tail two use $a\le m^2/2$.
Also $2B'\le a-1$ and $a(8B'-2)>6$, since $a\ge m^2/4$
and $m>10^{15}$.

Suppose $t\ge2$ for now.
Every genuine component-fixing entry costs at least $B'/a$ in
normalized moving rank, including entries whose natural layer
image is identity: apply (NMT.4) to the entire partner, or use
full rank for a genuine total scalar.  If its sum of local
natural indices is $I_i$, the same averaging over active
components gives the stronger lower bound
$AI_i/(ta)+B'/a$.

A component-moving entry has normalized support at least
$(a-1)/(2a)$.  This is the accepted exact weighted-cycle
word-orbit bound, with all other cycles and the entire partner
retained by convolution.  Four such entries together with
eight other genuine entries would cost
$2e+(8B'-2)e/a>2e+6$.  Thus
\[
               r_{\rm mov}\le3,\qquad r_{\rm fix}\ge9.
 \tag{NMT.13}
\]

Use the first, unshortened inequality of (TWP.10):
\[
 \sum_i\operatorname{rank}(1-T_i)-2e
 \ge\frac ea\left[r_{\rm fix}B'-2-\frac2t(L-1)\right].
 \tag{NMT.14}
\]
For clarity, this is not an invocation of the final HD-TW04
sufficient condition $a(9B'-L-1)>6$, which need not hold.
The preceding proof of (TWP.10) uses only the verified local
bound, $\lambda\ge0$, $A\ge0$, the whole partner floor,
and the actual quotient inequalities on $t$ and $mt$ letters.
The exact weighted-cycle loss is at most $1/a$ per nontrivial
cycle, and gives
$AI_i/(ta)+\lambda\operatorname{ind}(\pi_i)/t$ before summation.
Riemann--Hurwitz gives total indices at least $2mt-2$ and
$2t-2$, respectively.  This proves (NMT.14) with all original
identity projections and lift orders retained.

If $t\ge240$, (NMT.10)--(NMT.14) imply
\[
 \begin{split}
 \sum_i\operatorname{rank}(1-T_i)-2e
 &\ge\frac ea\left[\frac{a}{12m}-2-\frac{10a}{tm}\right]\\
 &\ge\frac ea\left[\frac{a}{24m}-2\right]
 \ge\frac{e}{48m}>6.
 \end{split}
 \tag{NMT.15}
\]
The penultimate step uses $a\ge96m$, and the last uses
$e\ge a^2\ge m^4/16$.
This contradicts (NMT.2).  Hence every surviving $t\ge2$
would satisfy
\[
                              2\le t\le239.
 \tag{NMT.16}
\]

\subsection*{3. Tails three through 240}

We record the exact accepted local comparisons.  Let $f(y)$
be the number of fixed natural letters.  In the first row,
$F_3(y)$ counts fixed unordered three-sets and
$B_3=\binom m3$.  The remaining rows compare with the ordered
cube fixed-point count.  All errors are localized to $f<m/2$.
\[
 \begin{array}{c|c|c|c}
 k&\text{comparator }\phi(y)&\epsilon_0&\text{local hypotheses}\\ \hline
 3&F_3(y)/B_3&48/m^{3/2}&m\ge64\\
 4&(f(y)/m)^3&4M_\mu^2/m^2&m\ge256M_\mu\\
 5,6&(f(y)/m)^3&2M_\mu(M_\mu+1)/m^3&m\ge256M_\mu\\
 7,\ldots,240&(f(y)/m)^3&2M_\mu/m^4&m\ge2^{k+1}M_\mu
 \end{array}
 \tag{NMT.17}
\]
These are HD-FT11, HD-FT09, HD-FT07 and HD-FT04, respectively.
They assert $|\chi(y)|/a\le\phi(y)$ when $f\ge m/2$ and
$|\chi(y)|/a\le\phi(y)+\epsilon_0$ otherwise, for every
permutation and every power.  Transposition only changes sign.
They also give $a\ge m^3/12$ in the first row and $a\ge2m^3$
in the other rows.

Every displayed local threshold holds here.  For $k\le6$,
(NMT.7) gives $M_\mu\le400$ at $k=4$, and $M_\mu<50000$
at $k=5,6$, so (NMT.5) suffices.
For $k\le240$,
$2^{k+1}M_\mu\le2^{241}(1+240!2^{240})
<4\cdot960^{240}<241^{482}<m$.
Also $M_\mu<m$.

Call an original entry bad if it moves components, or if it
fixes all components and some local permutation fixes fewer
than $m/2$ natural letters.  A nonbad entry has only nonbad
powers: a power fixes every letter fixed by the original
permutation.  Every bad entry has normalized support greater
than $1/5$.  For a moving entry use $(a-1)/(2a)$.
For a fixed entry use the local bad-entry estimate in the
proofs of HD-FT05, FT08, FT10 or FT12 and convolution.
For the cube rows it is
$\nu(W(y))/a>1/4-\epsilon_0>1/5$.
For tail three it is
$\nu(W(y))/a>1/4-7/m-48/m^{3/2}>1/5$.
The bounds hold at (NMT.5) and the thresholds just verified.
Since $e\ge a^2$ and $2+6/e<11/5$, at most ten original
entries are bad.

Apply the trace averaging (ICL.6)--(ICL.8) to the local
nonnegative comparator $\phi$.  If at least one component
is fixed, it gives the average of its fixed-component
comparators, with error at most $\epsilon_0$.
If none is fixed, the trace is at most $a^{-t/2}\le1/a$.
Thus the corresponding induced permutation action has error
\[
                 \epsilon=\max\{\epsilon_0,1/a\}
 \tag{NMT.18}
\]
at bad original entries, and zero error at every power of
a nonbad entry.  This comparison uses every original
coefficient order and arbitrary whole partner operators,
exactly as in D-TD49.  No new branch tuple is postulated.

For $k\ge4$ use the disjoint union of the $t$ ordered cubes.
The geometric part (ICL.14)--(ICL.17), which does not use
the special local spin/ordinary inputs of that theorem,
applies to every component-base genus.  With (NMT.16),
\[
 \frac1e\sum_i\operatorname{rank}(1-T_i)
 >2+\frac1{5736m}-\frac{10}{m^3}-10\epsilon.
 \tag{NMT.19}
\]
For $k=4$, one has $\epsilon=\epsilon_0\le640000/m^2$.
For $k=5,6$, one has $\epsilon=\epsilon_0<5100000000/m^3$.
For $k\ge7$, (NMT.17) and $M_\mu<m$ give
$\epsilon\le2/m^3$.
In all these cases the positive excess in (NMT.19) is greater
than $1/m^2$ at $m>10^{15}$.
For example, the worst quadratic-error case reduces to
$m/5736>6400001+10/m$; the other two cases reduce to weaker
inequalities at that same endpoint.
Since $6/e\le96/m^4<1/m^2$, this contradicts (NMT.2).

For $k=3$, use instead the disjoint union of the $t$ unordered
three-set actions.  It is one transitive actual quotient
action, of degree $tB_3$; let its curve be $Z_3$.
Let $X$ be the component-stabilizer curve and $h=g(X)$.
The actual $A_m/S_m$ component quotient over $X$ is the one
constructed in D-TD49.  If $h=0$, HD-GC06 gives
$g(Z_3)>m^2/288$.  If $h\ge2$, Riemann--Hurwitz gives
$g(Z_3)\ge1+B_3$, a stronger bound.
If $h=1$, nonabelian monodromy forces a branch $y\ne1$.
Choose distinct letters $u,v$ in one cycle of $y$.
A $y$-invariant subset contains $u$ if and only if it
contains $v$.  Therefore the $2\binom{m-2}{2}$ three-sets
containing exactly one of them are moved, and
$\operatorname{ind}_{\binom{[m]}3}(y)\ge\binom{m-2}{2}$.
It follows that $g(Z_3)\ge1+\binom{m-2}{2}/2>m^2/288$.
Thus the same genus lower bound holds in every case.

There are at most ten bad original entries, and
$1/a\le12/m^3\le48/m^{3/2}$.  The actual quotient identity
and the full-order Fourier comparison consequently give
\[
 \begin{split}
 \frac1e\sum_i\operatorname{rank}(1-T_i)
 &\ge2+\frac{2g(Z_3)-2}{tB_3}-\frac{480}{m^{3/2}}\\
 &>2+\frac1{24tm}-\frac{12}{tm^3}-\frac{480}{m^{3/2}}\\
 &\ge2+\frac1{5736m}-\frac6{m^3}-\frac{480}{m^{3/2}}
 >2+\frac1{m^2}>2+\frac6e .
 \end{split}
 \tag{NMT.20}
\]
For the penultimate inequality it is enough to use
$\sqrt m>10^7$ and
$1/5736>480/10^7+10^{-15}+6\cdot10^{-30}$.
These are exact rational endpoint comparisons.
This excludes all tails $3\le k\le240$ for $t\ge2$.

\subsection*{4. Tail two: the induced pair curve}

Suppose $k=2$ and $2\le t\le239$.  Write
$B_2=\binom m2$ and $F_2(y)$ for the number of fixed unordered
pairs.  The accepted HD-FT13 comparison is
\[
 a\ge m^2/4,\qquad
 \frac{|\chi(y)|}{a}\le\frac{F_2(y)}{B_2}
                  +\frac8{m^2}\,\mathbf1_{f(y)\le2}.
 \tag{NMT.21}
\]
It holds for both two-box labels, their transposes and all
powers.  A local element with $f(y)\le2$ has projective
support greater than $9a/20$, as proved in (FT2P.12).

Now call an original entry bad if it moves components or
fixes them all with some local $f\le2$.
Moving entries also have support greater than $9e/20$,
because $(a-1)/(2a)>9/20$ for $a>10$.
Hence there are at most four bad original entries, since
$2+6/e<9/4$.  Nonbad entries have only nonbad powers.
The induced action on the disjoint union of the within-block
unordered pairs is transitive and has degree $tB_2$.
The same trace average gives error at most $8/m^2$ at each
bad entry: if no component is fixed, use
$a^{-t/2}\le1/a\le4/m^2$.
Let $Z_2$ be this actual connected quotient curve, with genus
$\beta$.  The exact quotient identity gives
\[
 2+\frac6e\ge
 2+\frac{2\beta-2}{tB_2}-\frac{32}{m^2},
 \qquad
                  \beta<9t\le2151.
 \tag{NMT.22}
\]
Indeed $6/e\le96/m^4<1/m^2$ and $B_2<m^2/2$ give
$2\beta-2<(33/2)t$, while $t\ge2$.

Let $X$ again be the actual component-stabilizer curve.
If $g(X)\ge2$, then $\beta\ge1+B_2>2151$.
If $g(X)=1$, the actual nonabelian component cover has a
nontrivial branch $y$.  Choose two distinct letters in
one of its cycles.  The $2(m-2)$ unordered pairs containing
exactly one of them are moved, so the local pair index is
at least $m-2$ and $\beta\ge m/2>2151$.
Both cases contradict (NMT.22).  Therefore $X$ is a sphere.

\subsection*{5. Eight component branches, but twelve original marks}

The component cover over $X$ is now an actual $A_m/S_m$
sphere cover with unordered-pair genus $\beta\le2151$.
Let $g$ and $h$ be its natural and ordered-pair quotient
genera.  The same permutation-module inclusion used in
HD-GC07 gives $g\le\beta$.
Each even natural cycle contributes one to the ramification
of the ordered-to-unordered pair double cover, and at least
one to the natural index.  Thus the exact HD-GC01 formulas
imply
\[
 \begin{gathered}
 h\le m+2\beta+g-2,\\
 A=2(h-1)-2(m-1)(g-1)
       \le4m+4\beta-8-2(m-2)g
       \le4m+4\beta-8<6m .
 \end{gathered}
 \tag{NMT.23}
\]
Here $A$ is the total nonnegative relative ramification in
the ordered-pair-to-natural map; it is a sum of the local
cycle quantities in HD-GC03.

Apply the already proved $6m$ version of that cycle-mode
lemma.  Finite modes have exceptional length at most 252;
mode one moves at most twelve letters; mode infinity has
natural index at least $6m/7-72$.
Every large branch therefore has index at least $m/2-252$.
Five large branches would have index at least $5m/2-1260$,
greater than the actual total $2m+2g-2\le2m+4300$ at
$m>11120$.  Hence at most four branches are large.
If $J_0$ is the sum of the indices of the small branches,
their fixed-to-moved cycle contributions give
\[
 (m-12)J_0\le A\le4m+4\beta-8,\qquad J_0<5
                          \quad(m>8656).
 \tag{NMT.24}
\]
There are at most four small genuine branches.  The actual
component cover therefore has at most eight genuine branch
points.  No classification of their exact ramification types,
and no unevaluated pair-genus cutoff, has been used.

There are at most three original component-moving entries
by (NMT.13).  If an original entry fixes all components but
has a nonidentity natural layer image, the degree-$t$ map
$X\to\mathbf P^1$ is unramified above that branch value.
At least one of the $t$ points above it has nonidentity
local image in the component $A_m/S_m$ cover.
Distinct original branch values give distinct such points.
Thus at most eight additional original entries can fix
components but have nonidentity natural layer image.
At most eleven original entries have nonidentity natural
layer image in total.

Since $r\ge12$, some genuine original entry has identity
natural layer image.  Such a layer operator is scalar:
perfectness kills the discrepancy of its action on each
represented core, and irreducibility of the whole core
kills its remaining centralizer.  Its full actual operator
therefore has normalized moving rank at least $1/(108m)$
by (NMT.4), or full rank if it is a nonidentity total scalar.

On the other hand, sum the induced-pair Fourier lower
bound only over entries with nonidentity natural layer
image.  All quotient indices are retained, and at most
four errors occur.  Since $\beta\ge0$, their total is at
least $2-2/(tB_2)-32/m^2$.
The whole budget (NMT.2) leaves for all layer-identity
entries at most
\[
 \frac6e+\frac2{tB_2}+\frac{32}{m^2}
 <\frac{36}{m^2}<\frac1{108m}.
 \tag{NMT.25}
\]
We used $6/e<1/m^2$, $t\ge2$, $1/B_2\le3/m^2$ for
$m\ge3$, and $m>3888$.
This contradicts the preceding genuine-entry lower bound
and excludes tail two with $t\ge2$.

\subsection*{6. Singletons and the precise remaining layer type}

For $t=1$ and $k\ge3$, the accepted HD-FT17 effective
factorwise theorem already excludes the source at
$m\ge241^{482}$, with an arbitrary entire partner.
For $t=1$ and $k=2$, HD-FT14 and the effective HD-GC07
branch bound give at least four genuine original entries
acting only through the entire partner, with total
normalized moving rank at most $61/m^2$.
By (NMT.4) each costs at least $1/(108m)$, so
\[
               \frac4{108m}\le\frac{61}{m^2},
                         \qquad m\le1647,
 \tag{NMT.26}
\]
again impossible.  This is the already proved entire-partner
floor applied in the near-maximum range, not a new supplier
budget for $R$.

All cases $2\le k\le240$ and $t\ge1$ are excluded.
D-TD49 separately eliminates near-maximum spin, split and
ordinary components with degree at least $m^{120}$.
Its arbitrary-shape degree conclusion leaves precisely the
small ordinary tails considered here and the standard label.
A nontrivial surviving near-maximum alternating layer is
therefore ordinary standard.  This theorem leaves open
whether such a standard layer has one or more permuted
components.  It constructs no curve source and no cycle.
\end{proof}

\paragraph{Known inputs and audit boundary.}
All local character estimates, the stable-character coefficient
formula, ordinary branching and the tableau bound are exactly
the accepted inputs of HD-TW02--05 and HD-FT02/04/07/09/11/13/17.
The published character-polynomial and tableau sources retain
their already reviewed scopes.\cite{HD:GarsiaGoupil2009,HD:MSP2007}
The coefficient refinement uses only the sum of the same
partition-count bounds already proved there.
The cube and induced-set trace mechanism is D-TD49, with its
actual full-power normalization.  The sphere genus and cycle-mode
inputs are the fully proved HD-GC01/03/06/07.
The new proof above treats every positive component-base genus,
the bounded pair-genus version of the branch argument, and the
original-mark count explicitly.  It invokes no new external
classification theorem or unspecified asymptotic constant.
